\paragraph{A measurement belongs to a particular book.}
The optional \texttt{MixedEventTimeline/3} adds deferred assessment and an
explicit terminal carry to the preceding timeline. Deferred assessment is
dated by \code{measurement.at}; its supplied state digest must match the
financial book when the operation is reached. The runner retains those
measurement bytes. It does not replace the digest merely because an earlier
action has changed the book.

This matters even for a payment that creates no expense. Start with the
sale-first example, now assessed through 31~December 2027. Cash is USD~467,
profit before tax is USD~5 and current tax is USD~1.25. A deferred measurement
prepared for that state cannot be reused after paying USD~1 of a carried tax
liability: cash and payable have both changed. The original request stops at
the stale measurement and retains the payment. An author can instead supply
a distinct measurement for the resulting book, including its own identity,
date, tax bases and recovery assumptions. That is a new request, not an
automatic renewal of the old evidence.

The constructed deferred calculation also tests a nonzero accounting effect.
It begins with a deferred-tax asset and liability of USD~5 each. The supplied
taxable temporary difference of USD~20 keeps the liability at USD~5 under
the one-quarter rate. For the asset, the author now supplies only USD~12 of
recognized deductions and matching recovery capacity, giving an asset of
USD~3. Reducing the asset from USD~5 to USD~3 records USD~2 of deferred-tax
expense without changing cash. After-tax profit is consequently
USD~$5-1.25-2=1.75$. The recovery assumption is hypothetical; the example
does not estimate a company's future taxable income.

Carry occurs at the next period's start, which must equal the current period's
exclusive end. It requires current assessment through the closing day and a
fresh deferred assessment of the same book. In the case without the intervening
payment, carry on 1~January 2028 preserves cash of USD~467, inventory of USD~70
and current-tax payable of USD~11.50. It transfers profit of USD~1.75 to retained
earnings, whose credit-signed balance changes from USD~$-11.75$ to
USD~$-13.50$, and resets the new period's income and cash-flow movements.
The old statements, tax vintages and assessment history remain in the period
archive. The separately reported current-period profit is therefore zero
immediately after carry; it does not replace the archived profit of USD~1.75.

This bounded calendar permits one carry, as its last action. Later decisions
use another request bound to the new opening. Applying the retained inventory
rule on 2~January leaves its half-unit reserve intact: one unit remains before
the decision, so only half a unit is sold. Proceeds are USD~7.50 and cost is
USD~5, leaving cash of USD~474.50 and total inventory of USD~65. The preceding
year's report and tax vintages remain unchanged.

Thirteen focused tests and 61 regressions retain twelve complete exchanges:
\begin{flushleft}\small
\path{data/accounting-master/fiscal-timeline-20261025/}
\end{flushleft}
They compare entire native transitions, journals and archives, including an expressly renewed
measurement and a payment after deferred assessment that prevents carry.
An unknown deferred finding remains unresolved; an absent required entry is
invalid, and a known false prerequisite is unsupported. Provisional carry
retains its own unknown or false finding statuses without granting accounting
admission. Exact clocks and valid prefixes are preserved under the existing
32-native/64-total work limits. These examples specify timing and composition;
complete fiscal activity, general multi-period calendars and actual-company
tax findings remain separate obligations.
